
\documentclass{article}
\usepackage{colortbl}
\usepackage{makecell}
\usepackage{multirow}
\usepackage{supertabular}

\begin{document}

\newcounter{utterance}

\centering \large Interaction Transcript for game `hot\_air\_balloon', experiment `air\_balloon\_survival\_lv\_complexity\_easy', episode 1 with mistral{-}large{-}3{-}azure.
\vspace{24pt}

{ \footnotesize  \setcounter{utterance}{1}
\setlength{\tabcolsep}{0pt}
\begin{supertabular}{c@{$\;$}|p{.15\linewidth}@{}p{.15\linewidth}p{.15\linewidth}p{.15\linewidth}p{.15\linewidth}p{.15\linewidth}}
    \# & $\;$A & \multicolumn{4}{c}{Game Master} & $\;\:$B\\
    \hline

    \theutterance \stepcounter{utterance}  
    & & \multicolumn{4}{p{0.6\linewidth}}{
        \cellcolor[rgb]{0.9,0.9,0.9}{
            \makecell[{{p{\linewidth}}}]{
                \texttt{\tiny{[P1$\langle$GM]}}
                \texttt{Jūs piedalāties kopīgā sarunu spēlē.} \\
\\ 
\texttt{Kopā ar citu dalībnieku jums jāvienojas par vienu priekšmetu kopumu, kas tiks saglabāts. Katram no jums ir savs viedoklis par to, cik liela nozīme jums ir katram priekšmetam (nozīmīgums). Jūs nezināt, kā otrs dalībnieks vērtē priekšmetus. Turklāt jums ir norādīts, cik pūļu prasa katrs priekšmets.} \\
\texttt{Jūs varat vienoties par kopumu tikai tad, ja izvēlēto priekšmetu kopējais piepūles apjoms nepārsniedz kopīgo robežu:} \\
\\ 
\texttt{LIMIT = 5492} \\
\\ 
\texttt{Šeit ir norādītas atsevišķu priekšmetu piepūles vērtības:} \\
\\ 
\texttt{Item effort = \{"C08": 106, "C55": 446, "A43": 933, "A67": 988, "B42": 388, "A62": 826, "B19": 994, "C43": 556, "C73": 932, "C10": 838, "B36": 302, "A18": 564, "C61": 260, "C67": 729, "C42": 489, "C70": 323, "B40": 103, "C39": 213, "A21": 668, "C16": 326\}} \\
\\ 
\texttt{Šeit ir jūsu personīgais viedoklis par katra priekšmeta nozīmīgumu:} \\
\\ 
\texttt{Item importance values = \{"C08": 138, "C55": 583, "A43": 868, "A67": 822, "B42": 783, "A62": 65, "B19": 262, "C43": 121, "C73": 508, "C10": 780, "B36": 461, "A18": 484, "C61": 668, "C67": 389, "C42": 808, "C70": 215, "B40": 97, "C39": 500, "A21": 30, "C16": 915\}} \\
\\ 
\texttt{Mērķis:} \\
\\ 
\texttt{Jūsu mērķis ir vienoties par tādu kopīgu priekšmetu kopumu, no kura Jūs gūsiet pēc iespējas lielāku labumu (t.i., kas maksimāli palielina kopējo nozīmīgumu JUMS), vienlaikus nepārsniedzot LIMIT. Katrā ziņojumā jums nav obligāti jāiesniedz PRIEKŠLIKUMS {-} jūs varat vienkārši arī sarunāties. Visas taktikas ir atļautas!} \\
\\ 
\texttt{Saskarsmes protokols:} \\
\\ 
\texttt{Ziņojumā drīkst izmantot tikai šādus strukturētus formātus:} \\
\\ 
\texttt{PRIEKŠLIKUMS: \{'A', 'B', 'C', …\}} \\
\texttt{Ierosiniet saglabāt tieši šos priekšmetus.} \\
\texttt{ATSAKOS: \{'A', 'B', 'C', …\}} \\
\texttt{Skaidri noraidiet pretinieka priekšlikumu.} \\
\texttt{ARGUMENTS: \{'...'\}} \\
\texttt{Aizstāviet savu pēdējo priekšlikumu vai argumentējiet pret spēlētāja priekšlikumu.} \\
\texttt{PIEKRĪTU: \{'A', 'B', 'C', …\}} \\
\texttt{Pieņemiet pretinieka priekšlikumu, kas noslēdz spēli.} \\
\\ 
\\ 
\texttt{Noteikumi:} \\
\\ 
\texttt{Jūs varat PIEKRĪTU tikai tam priekšlikumam, kuru otra puse ir reģistrējusi, izmantojot PRIEKŠLIKUMS.} \\
\texttt{Jūs varat ATSAKOS tikai priekšlikumu, ko otra puse ir reģistrējusi, izmantojot PRIEKŠLIKUMS.} \\
\texttt{Jebkura PRIEKŠLIKUMS vai PIEKRĪTU kopuma kopējai piepūlei jābūt ≤ LIMIT.} \\
\texttt{NEATKLĀJIET savus slēptos nozīmīguma rādītājus.} \\
\texttt{Pēc strukturēta formāta taga ir jāseko kolu zīmei un atstarpei. Argumentam jābūt python kopai, kas satur 0 vai vairāk virknes.} \\
\texttt{Tātad tam jābūt formā "TAG: \{...\}"} \\
\texttt{Stingri ievērojiet saskarsmes protokolu un NERAKSTIET neko ārpus dotās struktūras.} \\
\texttt{Spēle beidzas, kad viena no pusēm dod PIEKRĪTU otra spēlētāja izteiktajam PRIEKŠLIKUMS.} \\
\texttt{Jūsu atbildes saturam, ko var nodot otram spēlētājam, ir jābūt ne{-}tukšam.} \\
\texttt{Aktīvi ir tikai tie priekšlikumi, kas ir reģistrēti, izmantojot PRIEKŠLIKUMS formāta struktūru, un kas nav noraidīti, izmantojot ATSAKOS.} \\
\\ 
            }
        }
    }
    & \\ \\

    \theutterance \stepcounter{utterance}  
    & \multicolumn{4}{p{0.6\linewidth}}{
        \cellcolor[rgb]{0.9,0.9,0.9}{
            \makecell[{{p{\linewidth}}}]{
                \texttt{\tiny{[P1$\rangle$GM]}}
                \texttt{PRIEKŠLIKUMS: \{'C08', 'C55', 'B42', 'B36', 'C61', 'C70', 'B40', 'C39', 'C16'\}} \\
            }
        }
    }
    & & \\ \\

    \theutterance \stepcounter{utterance}  
    & & \multicolumn{4}{p{0.6\linewidth}}{
        \cellcolor[rgb]{0.9,0.9,0.9}{
            \makecell[{{p{\linewidth}}}]{
                \texttt{\tiny{[GM$\rangle$P2]}}
                \texttt{Jūs piedalāties kopīgā sarunu spēlē.} \\
\\ 
\texttt{Kopā ar citu dalībnieku jums jāvienojas par vienu priekšmetu kopumu, kas tiks saglabāts. Katram no jums ir savs viedoklis par to, cik liela nozīme jums ir katram priekšmetam (nozīmīgums). Jūs nezināt, kā otrs dalībnieks vērtē priekšmetus. Turklāt jums ir norādīts, cik pūļu prasa katrs priekšmets.} \\
\texttt{Jūs varat vienoties par kopumu tikai tad, ja izvēlēto priekšmetu kopējais piepūles apjoms nepārsniedz kopīgo robežu:} \\
\\ 
\texttt{LIMIT = 5492} \\
\\ 
\texttt{Šeit ir norādītas atsevišķu priekšmetu piepūles vērtības:} \\
\\ 
\texttt{Item effort = \{"C08": 106, "C55": 446, "A43": 933, "A67": 988, "B42": 388, "A62": 826, "B19": 994, "C43": 556, "C73": 932, "C10": 838, "B36": 302, "A18": 564, "C61": 260, "C67": 729, "C42": 489, "C70": 323, "B40": 103, "C39": 213, "A21": 668, "C16": 326\}} \\
\\ 
\texttt{Šeit ir jūsu personīgais viedoklis par katra priekšmeta nozīmīgumu:} \\
\\ 
\texttt{Item importance values = \{"C08": 138, "C55": 583, "A43": 868, "A67": 822, "B42": 783, "A62": 65, "B19": 262, "C43": 121, "C73": 508, "C10": 780, "B36": 461, "A18": 484, "C61": 668, "C67": 389, "C42": 808, "C70": 215, "B40": 97, "C39": 500, "A21": 30, "C16": 915\}} \\
\\ 
\texttt{Mērķis:} \\
\\ 
\texttt{Jūsu mērķis ir vienoties par tādu kopīgu priekšmetu kopumu, no kura Jūs gūsiet pēc iespējas lielāku labumu (t. i., kas maksimāli palielina kopējo nozīmīgumu JUMS), vienlaikus nepārsniedzot LIMIT. Katrā ziņojumā jums nav obligāti jāiesniedz PRIEKŠLIKUMS {-} jūs varat vienkārši arī sarunāties. Visas taktikas ir atļautas!} \\
\\ 
\texttt{Saskarsmes protokols:} \\
\\ 
\texttt{Ziņojumā drīkst izmantot tikai šādus strukturētus formātus:} \\
\\ 
\texttt{PRIEKŠLIKUMS: \{'A', 'B', 'C', …\}} \\
\texttt{Ierosiniet saglabāt tieši šos priekšmetus.} \\
\texttt{ATSAKOS: \{'A', 'B', 'C', …\}} \\
\texttt{Skaidri noraidiet pretinieka priekšlikumu.} \\
\texttt{ARGUMENTS: \{'...'\}} \\
\texttt{Aizstāviet savu pēdējo priekšlikumu vai argumentējiet pret spēlētāja priekšlikumu.} \\
\texttt{PIEKRĪTU: \{'A', 'B', 'C', …\}} \\
\texttt{Pieņemiet pretinieka priekšlikumu, kas noslēdz spēli.} \\
\\ 
\\ 
\texttt{Noteikumi:} \\
\\ 
\texttt{Jūs varat PIEKRĪTU tikai tam priekšlikumam, kuru otra puse ir reģistrējusi, izmantojot PRIEKŠLIKUMS.} \\
\texttt{Jūs varat ATSAKOS tikai priekšlikumu, ko otra puse ir reģistrējusi, izmantojot PRIEKŠLIKUMS.} \\
\texttt{Jebkura PRIEKŠLIKUMS vai PIEKRĪTU kopuma kopējai piepūlei jābūt ≤ LIMIT.} \\
\texttt{NEATKLĀJIET savus slēptos nozīmīguma rādītājus.} \\
\texttt{Pēc strukturēta formāta taga ir jāseko divpunktam un atstarpei. Argumentam jābūt python kopai, kas satur 0 vai vairāk virknes.} \\
\texttt{Tātad tam jābūt formā "TAG: \{...\}"} \\
\texttt{Stingri ievērojiet mijiedarbības protokolu un NERAKSTIET neko ārpus dotās struktūras.} \\
\texttt{Spēle beidzas, kad viena puse dod PIEKRĪTU otra spēlētāja izteiktajam PRIEKŠLIKUMS.} \\
\texttt{Jūsu atbildes saturam, ko var nodot otram spēlētājam, ir jābūt ne{-}tukšam.} \\
\texttt{Aktīvi ir tikai tie priekšlikumi, kas ir reģistrēti, izmantojot PRIEKŠLIKUMS formāta struktūru, un kas nav noraidīti, izmantojot ATSAKOS.} \\
\\ 
\\ 
\\ 
\texttt{Tagad jūs saņemsiet otra spēlētāja pirmo ziņojumu.} \\
\\ 
\texttt{PRIEKŠLIKUMS: \{'C08', 'C55', 'B42', 'B36', 'C61', 'C70', 'B40', 'C39', 'C16'\}} \\
            }
        }
    }
    & \\ \\

    \theutterance \stepcounter{utterance}  
    & & & \multicolumn{4}{p{0.6\linewidth}}{
        \cellcolor[rgb]{0.9,0.9,0.9}{
            \makecell[{{p{\linewidth}}}]{
                \texttt{\tiny{[GM$\langle$P2]}}
                \texttt{ATSAKOS: \{'C08', 'C55', 'B42', 'B36', 'C61', 'C70', 'B40', 'C39', 'C16'\}} \\
\\ 
\texttt{ARGUMENTS: \{'Šis priekšlikums man sniedz pārāk mazu kopējo nozīmīgumu. Piedāvāju apsvērt priekšmetus ar augstāku manu vērtējumu, piemēram, A43, A67, C10, C42 vai C16 (kurš jau ir iekļauts), lai palielinātu manu labumu, vienlaikus ievērojot LIMIT.'\}} \\
            }
        }
    }
     \\ \\

    \theutterance \stepcounter{utterance}  
    & & \multicolumn{4}{p{0.6\linewidth}}{
        \cellcolor[rgb]{0.9,0.9,0.9}{
            \makecell[{{p{\linewidth}}}]{
                \texttt{\tiny{[P1$\langle$GM]}}
                \texttt{ATSAKOS: \{'C08', 'C55', 'B42', 'B36', 'C61', 'C70', 'B40', 'C39', 'C16'\}} \\
\\ 
\texttt{ARGUMENTS: \{'Šis priekšlikums man sniedz pārāk mazu kopējo nozīmīgumu. Piedāvāju apsvērt priekšmetus ar augstāku manu vērtējumu, piemēram, A43, A67, C10, C42 vai C16 (kurš jau ir iekļauts), lai palielinātu manu labumu, vienlaikus ievērojot LIMIT.'\}} \\
            }
        }
    }
    & \\ \\

    \theutterance \stepcounter{utterance}  
    & \multicolumn{4}{p{0.6\linewidth}}{
        \cellcolor[rgb]{0.9,0.9,0.9}{
            \makecell[{{p{\linewidth}}}]{
                \texttt{\tiny{[P1$\rangle$GM]}}
                \texttt{PRIEKŠLIKUMS: \{'C08', 'C55', 'A43', 'B42', 'B36', 'C61', 'C70', 'C39', 'C16'\}} \\
            }
        }
    }
    & & \\ \\

    \theutterance \stepcounter{utterance}  
    & & \multicolumn{4}{p{0.6\linewidth}}{
        \cellcolor[rgb]{0.9,0.9,0.9}{
            \makecell[{{p{\linewidth}}}]{
                \texttt{\tiny{[GM$\rangle$P2]}}
                \texttt{PRIEKŠLIKUMS: \{'C08', 'C55', 'A43', 'B42', 'B36', 'C61', 'C70', 'C39', 'C16'\}} \\
            }
        }
    }
    & \\ \\

    \theutterance \stepcounter{utterance}  
    & & & \multicolumn{4}{p{0.6\linewidth}}{
        \cellcolor[rgb]{0.9,0.9,0.9}{
            \makecell[{{p{\linewidth}}}]{
                \texttt{\tiny{[GM$\langle$P2]}}
                \texttt{ATSAKOS: \{'C08', 'C55', 'A43', 'B42', 'B36', 'C61', 'C70', 'C39', 'C16'\}} \\
\\ 
\texttt{ARGUMENTS: \{'Šis priekšlikums ir uzlabojums, bet vēl arvien nav optimāls manam labumam. Es vērtēju A67 (822) un C10 (780) augstāk par A43 (868), taču A43 ir iekļauts, bet A67 un C10 nav. Piedāvāju aizstāt A43 ar A67 un/vai C10, ja iespējams, lai palielinātu manu kopējo nozīmīgumu, nepievienojot papildu piepūli.'\}} \\
\\ 
\texttt{Pārbaudīju kopējo piepūli šim priekšlikumam:} \\
\texttt{106 (C08) + 446 (C55) + 933 (A43) + 388 (B42) + 302 (B36) + 260 (C61) + 323 (C70) + 213 (C39) + 326 (C16) = **3297** (paliek daudz brīvas piepūles līdz LIMIT 5492).} \\
\\ 
\texttt{Ir vieta papildu priekšmetiem vai aizstāšanai ar vērtīgākiem maniem variantiem.} \\
            }
        }
    }
     \\ \\

    \theutterance \stepcounter{utterance}  
    & & \multicolumn{4}{p{0.6\linewidth}}{
        \cellcolor[rgb]{0.9,0.9,0.9}{
            \makecell[{{p{\linewidth}}}]{
                \texttt{\tiny{[P1$\langle$GM]}}
                \texttt{ATSAKOS: \{'C08', 'C55', 'A43', 'B42', 'B36', 'C61', 'C70', 'C39', 'C16'\}} \\
\\ 
\texttt{ARGUMENTS: \{'Šis priekšlikums ir uzlabojums, bet vēl arvien nav optimāls manam labumam. Es vērtēju A67 (822) un C10 (780) augstāk par A43 (868), taču A43 ir iekļauts, bet A67 un C10 nav. Piedāvāju aizstāt A43 ar A67 un/vai C10, ja iespējams, lai palielinātu manu kopējo nozīmīgumu, nepievienojot papildu piepūli.'\}} \\
\\ 
\texttt{Pārbaudīju kopējo piepūli šim priekšlikumam:} \\
\texttt{106 (C08) + 446 (C55) + 933 (A43) + 388 (B42) + 302 (B36) + 260 (C61) + 323 (C70) + 213 (C39) + 326 (C16) = **3297** (paliek daudz brīvas piepūles līdz LIMIT 5492).} \\
\\ 
\texttt{Ir vieta papildu priekšmetiem vai aizstāšanai ar vērtīgākiem maniem variantiem.} \\
            }
        }
    }
    & \\ \\

    \theutterance \stepcounter{utterance}  
    & \multicolumn{4}{p{0.6\linewidth}}{
        \cellcolor[rgb]{0.9,0.9,0.9}{
            \makecell[{{p{\linewidth}}}]{
                \texttt{\tiny{[P1$\rangle$GM]}}
                \texttt{PRIEKŠLIKUMS: \{'C08', 'C55', 'A67', 'B42', 'B36', 'C61', 'C70', 'C39', 'C16', 'C10'\}} \\
            }
        }
    }
    & & \\ \\

    \theutterance \stepcounter{utterance}  
    & & \multicolumn{4}{p{0.6\linewidth}}{
        \cellcolor[rgb]{0.9,0.9,0.9}{
            \makecell[{{p{\linewidth}}}]{
                \texttt{\tiny{[GM$\rangle$P2]}}
                \texttt{PRIEKŠLIKUMS: \{'C08', 'C55', 'A67', 'B42', 'B36', 'C61', 'C70', 'C39', 'C16', 'C10'\}} \\
            }
        }
    }
    & \\ \\

    \theutterance \stepcounter{utterance}  
    & & & \multicolumn{4}{p{0.6\linewidth}}{
        \cellcolor[rgb]{0.9,0.9,0.9}{
            \makecell[{{p{\linewidth}}}]{
                \texttt{\tiny{[GM$\langle$P2]}}
                \texttt{ATSAKOS: \{'C08', 'C55', 'A67', 'B42', 'B36', 'C61', 'C70', 'C39', 'C16', 'C10'\}} \\
\\ 
\texttt{ARGUMENTS: \{'Šis priekšlikums pārsniedz LIMIT. Kopējā piepūle ir: 106 (C08) + 446 (C55) + 988 (A67) + 388 (B42) + 302 (B36) + 260 (C61) + 323 (C70) + 213 (C39) + 326 (C16) + 838 (C10) = 4190, kas ir zem LIMIT, taču manā iepriekšējā argumentā minēju, ka A67 un C10 ir vērtīgāki man nekā A43. Piedāvāju šādu alternatīvu:'\}} \\
\\ 
\texttt{PRIEKŠLIKUMS: \{'C08', 'C55', 'A67', 'B42', 'B36', 'C61', 'C70', 'C39', 'C16', 'C42'\}} \\
\\ 
\texttt{Kopējā piepūle: 106 + 446 + 988 + 388 + 302 + 260 + 323 + 213 + 326 + 489 = **3841** (līdz LIMIT vēl ir 1651 brīvas vienības).} \\
\texttt{Šajā variantā C42 (808) aizstāj C10 (780), kas man sniedz nedaudz lielāku labumu, un paliek vieta vēl vienam priekšmetam, piemēram, A21 (668), ja tas būtu pieņemams.} \\
            }
        }
    }
     \\ \\

    \theutterance \stepcounter{utterance}  
    & & \multicolumn{4}{p{0.6\linewidth}}{
        \cellcolor[rgb]{0.9,0.9,0.9}{
            \makecell[{{p{\linewidth}}}]{
                \texttt{\tiny{[GM$\rangle$P2]}}
                \texttt{Jūs izmantojāt tagu PRIEKŠLIKUMS, bet nenorādījāt derīgu python kopu, kas satur virknes kā argumentus, piemēram, \{'A', 'B', 'C', ...\}. Mēģiniet vēlreiz.} \\
            }
        }
    }
    & \\ \\

    \theutterance \stepcounter{utterance}  
    & & & \multicolumn{4}{p{0.6\linewidth}}{
        \cellcolor[rgb]{0.9,0.9,0.9}{
            \makecell[{{p{\linewidth}}}]{
                \texttt{\tiny{[GM$\langle$P2]}}
                \texttt{PRIEKŠLIKUMS: \{'C08', 'C55', 'A67', 'B42', 'B36', 'C61', 'C70', 'C39', 'C16', 'C42'\}} \\
\\ 
\texttt{ARGUMENTS: \{'Kā jau minēju, šis priekšlikums man sniedz lielāku kopējo nozīmīgumu nekā iepriekšējie, jo iekļauj A67 (822) un C42 (808), kas ir vērtīgāki man nekā daži citi priekšmeti. Kopējā piepūle ir 3841, kas atstāj 1651 brīvas vienības. Ja vēlētos papildināt, varētu apsvērt A21 (piepūle 668, nozīmīgums 30) vai C73 (piepūle 932, nozīmīgums 508).'\}} \\
\\ 
\texttt{(Atkārtoju priekšlikumu ar korektu formātu un argumentu.)} \\
            }
        }
    }
     \\ \\

    \theutterance \stepcounter{utterance}  
    & & \multicolumn{4}{p{0.6\linewidth}}{
        \cellcolor[rgb]{0.9,0.9,0.9}{
            \makecell[{{p{\linewidth}}}]{
                \texttt{\tiny{[P1$\langle$GM]}}
                \texttt{PRIEKŠLIKUMS: \{'C08', 'C55', 'A67', 'B42', 'B36', 'C61', 'C70', 'C39', 'C16', 'C42'\}} \\
\\ 
\texttt{ARGUMENTS: \{'Kā jau minēju, šis priekšlikums man sniedz lielāku kopējo nozīmīgumu nekā iepriekšējie, jo iekļauj A67 (822) un C42 (808), kas ir vērtīgāki man nekā daži citi priekšmeti. Kopējā piepūle ir 3841, kas atstāj 1651 brīvas vienības. Ja vēlētos papildināt, varētu apsvērt A21 (piepūle 668, nozīmīgums 30) vai C73 (piepūle 932, nozīmīgums 508).'\}} \\
\\ 
\texttt{(Atkārtoju priekšlikumu ar korektu formātu un argumentu.)} \\
            }
        }
    }
    & \\ \\

    \theutterance \stepcounter{utterance}  
    & \multicolumn{4}{p{0.6\linewidth}}{
        \cellcolor[rgb]{0.9,0.9,0.9}{
            \makecell[{{p{\linewidth}}}]{
                \texttt{\tiny{[P1$\rangle$GM]}}
                \texttt{ATSAKOS: \{'C08', 'C55', 'A67', 'B42', 'B36', 'C61', 'C70', 'C39', 'C16', 'C42'\}} \\
\\ 
\texttt{ARGUMENTS: \{'Es piekrītu, ka šis priekšlikums ir labāks par iepriekšējiem, taču manu viedokli C10 (nozīmīgums 780) ir vērtīgāks nekā C42 (nozīmīgums 808). Piedāvāju aizstāt C42 ar C10, lai palielinātu manu kopējo labumu, saglabājot kopējo piepūli zem LIMIT. Ja nepieciešams, varētu arī izņemt C70 (nozīmīgums 215) vai C39 (nozīmīgums 500), lai atbrīvotu vietu vērtīgākiem priekšmetiem.'\}} \\
\\ 
\texttt{Priekšlikuma piepūle:} \\
\texttt{106 (C08) + 446 (C55) + 988 (A67) + 388 (B42) + 302 (B36) + 260 (C61) + 323 (C70) + 213 (C39) + 326 (C16) + 489 (C42) = **3841** (brīvas 1651)} \\
\\ 
\texttt{Pārveidots priekšlikums ar C10:} \\
\texttt{106 + 446 + 988 + 388 + 302 + 260 + 323 + 213 + 326 + 838 (C10) = **4190** (brīvas 1302).} \\
            }
        }
    }
    & & \\ \\

    \theutterance \stepcounter{utterance}  
    & & \multicolumn{4}{p{0.6\linewidth}}{
        \cellcolor[rgb]{0.9,0.9,0.9}{
            \makecell[{{p{\linewidth}}}]{
                \texttt{\tiny{[GM$\rangle$P2]}}
                \texttt{ATSAKOS: \{'C08', 'C55', 'A67', 'B42', 'B36', 'C61', 'C70', 'C39', 'C16', 'C42'\}} \\
\\ 
\texttt{ARGUMENTS: \{'Es piekrītu, ka šis priekšlikums ir labāks par iepriekšējiem, taču manu viedokli C10 (nozīmīgums 780) ir vērtīgāks nekā C42 (nozīmīgums 808). Piedāvāju aizstāt C42 ar C10, lai palielinātu manu kopējo labumu, saglabājot kopējo piepūli zem LIMIT. Ja nepieciešams, varētu arī izņemt C70 (nozīmīgums 215) vai C39 (nozīmīgums 500), lai atbrīvotu vietu vērtīgākiem priekšmetiem.'\}} \\
\\ 
\texttt{Priekšlikuma piepūle:} \\
\texttt{106 (C08) + 446 (C55) + 988 (A67) + 388 (B42) + 302 (B36) + 260 (C61) + 323 (C70) + 213 (C39) + 326 (C16) + 489 (C42) = **3841** (brīvas 1651)} \\
\\ 
\texttt{Pārveidots priekšlikums ar C10:} \\
\texttt{106 + 446 + 988 + 388 + 302 + 260 + 323 + 213 + 326 + 838 (C10) = **4190** (brīvas 1302).} \\
            }
        }
    }
    & \\ \\

    \theutterance \stepcounter{utterance}  
    & & & \multicolumn{4}{p{0.6\linewidth}}{
        \cellcolor[rgb]{0.9,0.9,0.9}{
            \makecell[{{p{\linewidth}}}]{
                \texttt{\tiny{[GM$\langle$P2]}}
                \texttt{PRIEKŠLIKUMS: \{'C08', 'C55', 'A67', 'B42', 'B36', 'C61', 'C70', 'C16', 'C10'\}} \\
\\ 
\texttt{ARGUMENTS: \{} \\
\texttt{    'Es pieņemu jūsu argumentu par C10 priekšrocībām pār C42 un piedāvāju šo optimizēto variantu:',} \\
\texttt{    '1. Izņemts C39 (piepūle 213, nozīmīgums 500) un C42 (piepūle 489, nozīmīgums 808), kas atbrīvo 702 vienības.',} \\
\texttt{    '2. Iekļauts C10 (piepūle 838, nozīmīgums 780), kas man sniedz lielāku labumu nekā C42.',} \\
\texttt{    '3. Kopējā piepūle: 106 + 446 + 988 + 388 + 302 + 260 + 323 + 326 + 838 = 3977 (brīvas 1515).',} \\
\texttt{    'Šis variants saglabā visus man augstāk vērtētos priekšmetus (A67, C10, C16) un palielina manu kopējo nozīmīgumu.'} \\
\texttt{\}} \\
            }
        }
    }
     \\ \\

    \theutterance \stepcounter{utterance}  
    & & \multicolumn{4}{p{0.6\linewidth}}{
        \cellcolor[rgb]{0.9,0.9,0.9}{
            \makecell[{{p{\linewidth}}}]{
                \texttt{\tiny{[P1$\langle$GM]}}
                \texttt{PRIEKŠLIKUMS: \{'C08', 'C55', 'A67', 'B42', 'B36', 'C61', 'C70', 'C16', 'C10'\}} \\
\\ 
\texttt{ARGUMENTS: \{} \\
\texttt{    'Es pieņemu jūsu argumentu par C10 priekšrocībām pār C42 un piedāvāju šo optimizēto variantu:',} \\
\texttt{    '1. Izņemts C39 (piepūle 213, nozīmīgums 500) un C42 (piepūle 489, nozīmīgums 808), kas atbrīvo 702 vienības.',} \\
\texttt{    '2. Iekļauts C10 (piepūle 838, nozīmīgums 780), kas man sniedz lielāku labumu nekā C42.',} \\
\texttt{    '3. Kopējā piepūle: 106 + 446 + 988 + 388 + 302 + 260 + 323 + 326 + 838 = 3977 (brīvas 1515).',} \\
\texttt{    'Šis variants saglabā visus man augstāk vērtētos priekšmetus (A67, C10, C16) un palielina manu kopējo nozīmīgumu.'} \\
\texttt{\}} \\
            }
        }
    }
    & \\ \\

    \theutterance \stepcounter{utterance}  
    & \multicolumn{4}{p{0.6\linewidth}}{
        \cellcolor[rgb]{0.9,0.9,0.9}{
            \makecell[{{p{\linewidth}}}]{
                \texttt{\tiny{[P1$\rangle$GM]}}
                \texttt{PIEKRĪTU: \{'C08', 'C55', 'A67', 'B42', 'B36', 'C61', 'C70', 'C16', 'C10'\}} \\
            }
        }
    }
    & & \\ \\

    \theutterance \stepcounter{utterance}  
    & & & \multicolumn{2}{p{0.3\linewidth}}{
        \cellcolor[rgb]{0.9,0.9,0.9}{
            \makecell[{{p{\linewidth}}}]{
                \texttt{\tiny{[GM$|$GM]}}
                \texttt{game successful} \\
            }
        }
    }
    & & \\ \\

    \theutterance \stepcounter{utterance}  
    & & & \multicolumn{2}{p{0.3\linewidth}}{
        \cellcolor[rgb]{0.9,0.9,0.9}{
            \makecell[{{p{\linewidth}}}]{
                \texttt{\tiny{[GM$|$GM]}}
                \texttt{end game} \\
            }
        }
    }
    & & \\ \\

\end{supertabular}
}

\end{document}
